\documentclass[../../main.tex]{subfiles}

\begin{document}

\chapter{Landasan Teori dan Konsep Dasar}
\label{chap:landasan-teori}

% =========================================================
% OUTCOME PEMBELAJARAN KHUSUS BAB
% =========================================================

\begin{subcpmkbox}
Bab ini tidak menetapkan Sub-CPMK baru, melainkan menyediakan
fondasi konseptual yang menopang \SubCPMK{1.1} sampai
\SubCPMK{4.3}.

\vspace{0.4em}
\textbf{CPMK yang didukung:} \CPMK{01}, \CPMK{02},
\CPMK{03}, \CPMK{04} \quad\textbullet\quad
\textbf{CPL yang didukung:} \CPL{02}, \CPL{03}
\end{subcpmkbox}

\begin{tujuanbox}
Setelah mempelajari bab ini, mahasiswa mampu:
\begin{enumerate}
    \item menjelaskan model mesin von Neumann dan prinsip
          \eng{stored-program};
    \item menjelaskan pengertian \eng{instruction set architecture},
          \eng{operand}, \eng{mnemonic}, register, dan siklus
          instruksi;
    \item mengonversi bilangan antara biner, heksadesimal, dan
          desimal tanpa alat bantu;
    \item menghitung nilai bertanda dengan representasi
          \eng{two's complement} dan mendeteksi \eng{overflow};
    \item menjelaskan representasi BCD dan ASCII serta kaitannya
          dengan operasi masukan/keluaran;
    \item membaca peta konsep buku dan menjelaskan urutan
          ketergantungan antarbab.
\end{enumerate}
\end{tujuanbox}

\begin{materibox}
\begin{enumerate}
    \item Model mesin von Neumann dan prinsip \eng{stored-program}.
    \item Konsep dan definisi kunci: ISA, \eng{operand},
          \eng{mnemonic}, register, siklus instruksi, dan
          \eng{memory-mapped I/O}.
    \item Representasi data pada level mesin: biner, heksadesimal,
          \eng{two's complement}, BCD, dan ASCII.
    \item Peta konsep buku dan urutan ketergantungan antarbab.
\end{enumerate}
\end{materibox}

\subfile{section-01}
\subfile{section-02}
\subfile{section-03}

\section{Contoh Penerapan}
\label{sec:contoh-bab2}

\subsection{Contoh 1: Konversi Bilangan Bertahap}

Berikut contoh konversi yang dikerjakan tanpa alat bantu, dengan
langkah yang dapat diperiksa.

\begin{enumerate}
    \item \textbf{Desimal 250 ke biner.}
          Bagi berulang oleh 2 dan catat sisanya:
          $250 = 1111\,1010_2$.

    \item \textbf{Biner \bin{1111\,1010} ke heksadesimal.}
          Kelompokkan per empat bit dari kanan: \bin{1111} dan
          \bin{1010}, yaitu \hex{F} dan \hex{A}. Hasilnya
          \hex{FA}.

    \item \textbf{Heksadesimal \hex{FA} ke desimal.}
          $15 \times 16 + 10 = 240 + 10 = 250$. Hasilnya konsisten
          dengan langkah pertama, sehingga ketiga konversi saling
          memverifikasi.
\end{enumerate}

\subsection{Contoh 2: Menghitung \eng{Two's Complement}}

Untuk menyatakan \dec{-25} dalam 8 bit:

\begin{enumerate}
    \item Tulis nilai positifnya dalam 8 bit: \bin{0001\,1001}
          (yaitu \hex{19}).
    \item Balikkan seluruh bitnya: \bin{1110\,0110} (yaitu
          \hex{E6}).
    \item Tambahkan satu: \bin{1110\,0111} (yaitu \hex{E7}).
\end{enumerate}

Verifikasi: \hex{E7} + \hex{19} = \hex{100}, yang dalam 8 bit menjadi
\hex{00} dengan \eng{carry} keluar. Hasilnya nol, sebagaimana
mestinya untuk $(-25) + 25$. Sifat inilah yang membuat
\eng{two's complement} dipakai: prosesor tidak memerlukan rangkaian
khusus untuk pengurangan, cukup penjumlahan biasa.

\subsection{Contoh 3: Melihat Penyandian Data di Memori}

Program pada Kode~\ref{lst:penyandian} mendeklarasikan beberapa data
dengan penyandian berbeda, lalu mahasiswa memeriksa bagaimana
byte-byte tersebut benar-benar tersimpan.

\begin{lstlisting}[
    language=MASM,
    caption={Program yang memperlihatkan penyandian data di memori.},
    label={lst:penyandian}
]
.MODEL SMALL
.STACK 100h

.DATA
    nilai_desimal   DB  25          ; nilai biner, bukan sandian karakter
    nilai_bcd       DB  25h         ; 25 dalam BCD
    sandian_karakter DB '2'         ; sandian ASCII untuk karakter '2'
    kata_16bit      DW  1234h       ; dua byte, little-endian

.CODE
START:
    MOV AX, @DATA
    MOV DS, AX

    MOV AH, 4Ch
    MOV AL, 00h
    INT 21h

END START
\end{lstlisting}

Setelah program dirakit, periksa isi segmen data dengan debugger.
Tabel~\ref{tab:penyandian-hasil} memperlihatkan hasil yang
seharusnya Anda temukan.

\begin{table}[H]
    \centering
    \caption{Isi memori untuk data pada Kode~\ref{lst:penyandian}.}
    \label{tab:penyandian-hasil}
    \small
    \begin{tabular}{@{}llp{7.2cm}@{}}
        \toprule
        \textbf{Data} & \textbf{Isi memori} & \textbf{Penjelasan} \\
        \midrule
        \code{nilai\_desimal} & \hex{19} &
        \dec{25} dalam biner; jika dicetak langsung, yang muncul
        adalah karakter tak terlihat, bukan ``25'' \\
        \addlinespace
        \code{nilai\_bcd} & \hex{25} &
        Dua digit BCD dalam satu byte: \bin{0010} dan \bin{0101} \\
        \addlinespace
        \code{sandian\_karakter} & \hex{32} &
        Sandian ASCII karakter \code{'2'}, yaitu \hex{30} + 2 \\
        \addlinespace
        \code{kata\_16bit} & \hex{34} lalu \hex{12} &
        Nilai \hex{1234} disimpan \eng{little-endian}: byte rendah
        lebih dahulu \\
        \bottomrule
    \end{tabular}
\end{table}

\begin{infobox}
Perhatikan tiga baris pertama: ketiganya ``bernilai 25'' menurut
manusia, tetapi tiga-tiganya tersimpan berbeda di memori. Inilah
sebabnya mahasiswa pemrograman tingkat rendah harus selalu
memikirkan \emph{representasi}, bukan hanya \emph{nilai}.
\end{infobox}

\subsection{Contoh 4: Menjumlahkan Data dan Memeriksa Flag}

Program pada Kode~\ref{lst:flag-uji} memperlihatkan bagaimana
representasi yang sama menghasilkan tafsiran berbeda.

\begin{lstlisting}[
    language=MASM,
    caption={Penjumlahan yang menghasilkan \eng{carry} dan
    \eng{overflow} sekaligus.},
    label={lst:flag-uji}
]
.MODEL SMALL
.STACK 100h

.CODE
START:
    MOV AX, @DATA
    MOV DS, AX

    MOV AL, 0FFh        ; 255 tak bertanda, atau -1 bertanda
    ADD AL, 01h         ; 255 + 1

    MOV AH, 4Ch
    MOV AL, 00h
    INT 21h

END START
\end{lstlisting}

Setelah \asm{ADD AL, 01h} dieksekusi: \reg{AL} bernilai \hex{00},
\flag{CF} bernilai 1 (ada \eng{carry} keluar dari bit ke-8),
\flag{ZF} bernilai 1 (hasilnya nol), dan \flag{OF} bernilai 0
(tidak ada \eng{overflow} bertanda karena $-1 + 1 = 0$ masih dalam
rentang). Tafsiran \eng{carry} dan \eng{overflow} bergantung pada
apakah program bermaksud memperlakukan \hex{FF} sebagai
\dec{255} atau sebagai $-1$.

\begin{session}[caption={Pengamatan register dan flag pada
Kode~\ref{lst:flag-uji}.}]
-R
AX=00FF  BX=0000  CX=0000  DX=0000  SP=FFFE  BP=0000  SI=0000  DI=0000
DS=0700  ES=0700  SS=0700  CS=0700  IP=0106   NV UP EI PL NZ NA PO NC
-T
AX=0000  BX=0000  CX=0000  DX=0000  SP=FFFE  BP=0000  SI=0000  DI=0000
DS=0700  ES=0700  SS=0700  CS=0700  IP=0108   NV UP EI PL ZR AC PE CY
\end{session}

Baris pertama memperlihatkan keadaan sebelum \asm{ADD AL, 01h}:
\reg{AL} berisi \hex{FF} dan seluruh flag utama bernilai 0, yang pada
DEBUG ditandai \code{NV} (tidak \eng{overflow}), \code{NZ} (tidak
nol), dan \code{NC} (tidak ada \eng{carry}).

Baris kedua memperlihatkan keadaan sesudahnya. Penanda berubah
menjadi \code{ZR} (\flag{ZF} = 1, hasilnya nol), \code{CY}
(\flag{CF} = 1, ada \eng{carry} keluar bit ke-8), dan \code{AC}
(\flag{AF} = 1, ada \eng{carry} dari \eng{nibble} rendah ke tinggi).
Penanda \code{NV} tetap, yang berarti \flag{OF} = 0: sebagai
bilangan bertanda, $-1 + 1 = 0$ tidak melampaui rentang sehingga
tidak terjadi \eng{overflow}.

Inilah contoh paling ringkas mengapa \eng{carry} dan
\eng{overflow} harus dibedakan: pada operasi yang sama,
\flag{CF} bernilai 1 sementara \flag{OF} bernilai 0.

\section{Praktikum}
\label{sec:praktikum-bab2}

\begin{praktikumbox}
\textbf{Judul:} Konversi bilangan dan verifikasi penyandian data di
memori.\\[0.3em]
\textbf{Sub-CPMK:} fondasi untuk 1.1--4.3 \quad\textbullet\quad
\textbf{Alokasi waktu:} 100 menit \quad\textbullet\quad
\textbf{Bentuk:} individu
\end{praktikumbox}

\subsection{Tujuan Praktikum}

Memastikan mahasiswa dapat melakukan konversi bilangan secara manual
dan memverifikasi sendiri bahwa penyandian data di memori sesuai
dengan yang dihitung.

\subsection{Alat dan Bahan}

\begin{itemize}
    \item Komputer dengan Emu8086 dan DOSBox terpasang.
    \item Program pada Kode~\ref{lst:penyandian} dan
          Kode~\ref{lst:flag-uji}.
    \item Kalkulator \eng{programmer} (boleh kalkulator bawaan sistem
          operasi) untuk memeriksa hasil, bukan untuk menggantikan
          perhitungan manual.
\end{itemize}

\subsection{Langkah Kerja}

\begin{enumerate}
    \item \textbf{Konversi manual.}
    \begin{enumerate}
        \item Konversi desimal \dec{200}, \dec{100}, dan \dec{63} ke
              biner dan heksadesimal tanpa alat bantu.
        \item Nyatakan \dec{-1}, \dec{-64}, dan \dec{-100} dalam
              \eng{two's complement} 8 bit.
        \item Periksa hasil Anda dengan kalkulator \eng{programmer}.
              Bila berbeda, temukan letak kesalahannya.
    \end{enumerate}

    \item \textbf{Verifikasi penyandian di memori.}
    \begin{enumerate}
        \item Rakit dan jalankan program pada Kode~\ref{lst:penyandian}
              di Emu8086.
        \item Buka jendela memori dan arahkan ke segmen data.
        \item Isi tabel pengamatan: untuk setiap variabel, catat
              alamat dan isi byte-nya.
        \item Bandingkan hasil pengamatan dengan perhitungan manual
              Anda pada Tabel~\ref{tab:penyandian-hasil}.
    \end{enumerate}

    \item \textbf{Menguji pengaruh flag.}
    \begin{enumerate}
        \item Jalankan program pada Kode~\ref{lst:flag-uji} dengan
              \eng{single step}.
        \item Sebelum eksekusi \asm{ADD AL, 01h}, tuliskan prediksi
              Anda untuk \flag{CF}, \flag{ZF}, \flag{SF}, dan
              \flag{OF}.
        \item Setelah eksekusi, catat nilai sebenarnya dan bandingkan
              dengan prediksi.
        \item Ulangi dengan nilai \asm{7Fh} serta \asm{80h} dan
              jelaskan perbedaan hasilnya.
    \end{enumerate}

    \item \textbf{Menyusun tabel konversi pribadi.}
    \begin{enumerate}
        \item Buat tabel heksadesimal 0--F beserta nilai biner dan
              desimalnya di buku catatan Anda.
        \item Hafalkan tabel tersebut. Tabel ini akan sangat
              mempercepat pekerjaan Anda pada bab-bab berikutnya.
    \end{enumerate}
\end{enumerate}

\subsection{Hasil yang Dikumpulkan}

\begin{enumerate}
    \item Lembar kerja konversi manual (sebelum diperiksa dengan
          kalkulator).
    \item Tabel pengamatan alamat dan isi memori.
    \item Tabel prediksi versus hasil pengamatan flag.
    \item Tabel konversi heksadesimal pribadi.
\end{enumerate}

\begin{asesmenbox}
\textbf{Asesmen fondasi (prasyarat, tanpa nilai numerik).} Praktikum
ini tidak dinilai, tetapi menjadi prasyarat untuk Tugas 1 pada
Bab~\ref{chap:register}. Mahasiswa yang belum menguasai konversi
bilangan akan mengalami kesulitan besar pada bab-bab berikutnya.

\begin{tabular}{@{}p{7.0cm}p{6.4cm}@{}}
    \toprule
    \textbf{Butir} & \textbf{Kriteria terpenuhi} \\
    \midrule
    Konversi manual & Seluruh soal konversi benar dan langkah
    pengerjaan terlihat \\
    \addlinespace
    Verifikasi memori & Isi memori hasil pengamatan sama dengan
    perhitungan manual \\
    \addlinespace
    Prediksi flag & Minimal tiga dari empat prediksi benar \\
    \addlinespace
    Tabel pribadi & Tabel 0--F lengkap dan sudah dihafalkan \\
    \bottomrule
\end{tabular}
\end{asesmenbox}

\section{Latihan}
\label{sec:latihan-bab2}

\begin{enumerate}
    \item \textbf{Konversi dasar.} Konversi bilangan berikut tanpa
          alat bantu: (a) \dec{173} ke biner dan heksadesimal;
          (b) \hex{5C} ke biner dan desimal; (c) \bin{1101\,0011} ke
          heksadesimal dan desimal.

    \item \textbf{\eng{Two's complement}.} Nyatakan bilangan berikut
          dalam 8-bit \eng{two's complement}: (a) \dec{-1};
          (b) \dec{-16}; (c) \dec{-100}. Tunjukkan langkah pembalikan
          bit dan penambahan satu untuk setiap soal.

    \item \textbf{Rentang nilai.} Berapa nilai terbesar dan terkecil
          yang dapat dinyatakan dalam 8 bit secara bertanda? Bagaimana
          bila dinyatakan secara tak bertanda? Jelaskan mengapa
          rentangnya berbeda.

    \item \textbf{Analisis flag.} Diberikan \reg{AL} = \hex{7F} dan
          instruksi \asm{ADD AL, 01h}. Tentukan hasil di \reg{AL}
          serta nilai \flag{CF}, \flag{ZF}, \flag{SF}, dan
          \flag{OF}. Jelaskan mengapa \flag{OF} bernilai 1 sementara
          \flag{CF} bernilai 0.

    \item \textbf{Analisis flag.} Ulangi soal 4 dengan \reg{AL} =
          \hex{FF}. Tentukan hasil dan seluruh flag, lalu jelaskan
          mengapa kali ini \flag{CF} bernilai 1 tetapi \flag{OF}
          bernilai 0.

    \item \textbf{Penyandian.} Sebuah variabel dideklarasikan sebagai
          \code{DB '5'}. Berapa nilai byte yang tersimpan? Berapa
          nilainya bila dideklarasikan sebagai \code{DB 5}? Jelaskan
          perbedaan dan akibatnya bila variabel itu dicetak dengan
          layanan DOS fungsi \hex{02}.

    \item \textbf{\eng{Little-endian}.} Sebuah variabel
          dideklarasikan \code{DW 0ABCDh}. Tuliskan urutan byte yang
          tersimpan di memori dari alamat terendah ke tertinggi, lalu
          jelaskan mengapa urutannya demikian.

    \item \textbf{Kaitan ke bab berikutnya.} Berdasarkan
          Gambar~\ref{fig:peta-konsep}, jelaskan mengapa
          Bab~\ref{chap:segmentasi} harus dipelajari sebelum
          Bab~\ref{chap:mode-pengalamatan}. Kaitkan jawaban Anda
          dengan konsep alamat efektif.
\end{enumerate}

\begin{refleksibox}
Jawab pertanyaan berikut pada jurnal praktikum Anda.
\begin{enumerate}
    \item Bagian mana dari representasi data yang paling menantang
          bagi Anda: konversi bilangan, \eng{two's complement}, atau
          pembedaan \eng{carry} dan \eng{overflow}? Mengapa?
    \item Setelah melihat isi memori secara langsung, apakah
          pemahaman Anda tentang ``penyimpanan data'' berubah?
          Bagian mana yang paling mengejutkan?
    \item Konsep apa pada bab ini yang Anda rasa akan paling sering
          Anda pakai pada bab-bab berikutnya?
\end{enumerate}
\end{refleksibox}

\section{Rangkuman}
\label{sec:rangkuman-bab2}

Pada bab ini telah dibahas:

\begin{itemize}
    \item Model mesin von Neumann terdiri atas unit pemroses, memori,
          unit masukan, dan unit keluaran. Ciri utamanya adalah
          prinsip \eng{stored-program}: instruksi dan data disimpan
          di memori yang sama.

    \item Siklus instruksi terdiri atas \eng{fetch}, \eng{decode},
          dan \eng{execute}, dengan \eng{store} sebagai langkah
          tambahan pada 8086. Satu instruksi dapat memerlukan
          beberapa siklus mesin.

    \item ISA adalah kontrak instruksi antara perangkat lunak dan
          perangkat keras. Kompatibilitas ISA membuat program 8086
          masih dapat dijalankan pada prosesor x86 modern.

    \item \eng{Opcode} menyatakan operasi, \eng{operand} menyatakan
          data, dan \eng{mnemonic} adalah nama pendek opcode dalam
          bahasa assembly.

    \item Register adalah penyimpan kecil di dalam prosesor dengan
          akses jauh lebih cepat daripada memori. Jumlahnya terbatas,
          sehingga harus dikelola secara sadar.

    \item Penyandian data di level mesin meliputi biner,
          heksadesimal, \eng{two's complement}, BCD, dan ASCII.
          Representasi yang sama dapat memiliki tafsiran berbeda
          bergantung pada apakah bilangan diperlakukan bertanda atau
          tak bertanda.

    \item \eng{Carry} (\flag{CF}) menandakan kelebihan pada operasi
          tak bertanda, sedangkan \eng{overflow} (\flag{OF})
          menandakan kelebihan pada operasi bertanda. Keduanya harus
          dibedakan.

    \item Peta konsep buku menunjukkan empat klaster kompetensi yang
          saling bergantung, dan urutan bab mengikuti
          ketergantungan tersebut.
\end{itemize}

\begin{checklistbox}
Tandai setiap butir yang sudah Anda kuasai.

\begin{itemize}
    \item[\checkbox] Saya dapat menjelaskan model von Neumann dan
          prinsip \eng{stored-program}.
    \item[\checkbox] Saya dapat menjelaskan tahapan siklus instruksi.
    \item[\checkbox] Saya dapat mendefinisikan ISA, \eng{opcode},
          \eng{operand}, \eng{mnemonic}, dan register dengan kalimat
          saya sendiri.
    \item[\checkbox] Saya dapat mengonversi bilangan antara biner,
          heksadesimal, dan desimal tanpa alat bantu.
    \item[\checkbox] Saya dapat menyatakan bilangan negatif dalam
          \eng{two's complement} dan memverifikasinya.
    \item[\checkbox] Saya dapat membedakan \eng{carry} dan
          \eng{overflow} beserta flag yang menandainya.
    \item[\checkbox] Saya dapat menjelaskan penyandian BCD dan ASCII
          serta perbedaan nilai dengan sandian karakter.
    \item[\checkbox] Saya dapat membaca peta konsep buku dan
          menjelaskan urutan ketergantungan antarbab.
\end{itemize}

\vspace{0.4em}
Bila masih ada butir yang belum bertanda, ulangi bagian yang
bersesuaian sebelum melanjutkan ke Bab~\ref{chap:pengenalan}.
\end{checklistbox}

\begin{infobox}
\textbf{Rujukan bab.} Materi pada bab ini diperkuat oleh
\cite{tanenbaum2013}, \cite{stallings2016}, dan \cite{brey2009}
untuk organisasi komputer dan representasi data. Untuk gambaran
teknis 8086 yang lebih rinci, lihat \cite{intel1979}. Pembahasan
\eng{two's complement} dan operasi bertanda pada bahasa assembly
dapat dilihat pada \cite{hyde2010}.
\end{infobox}

\section{Evaluasi}
\label{sec:evaluasi-bab2}

\subsection{Soal Pilihan Ganda}

\begin{enumerate}
    \item Prinsip yang menyatakan bahwa instruksi dan data disimpan
          pada memori yang sama disebut \dots
          \begin{enumerate}
              \item \eng{pipelining}
              \item \eng{stored-program}
              \item \eng{memory-mapped I/O}
              \item \eng{little-endian}
          \end{enumerate}

    \item Urutan tahapan siklus instruksi yang benar adalah \dots
          \begin{enumerate}
              \item \eng{decode} $\rightarrow$ \eng{fetch}
                    $\rightarrow$ \eng{execute}
              \item \eng{execute} $\rightarrow$ \eng{fetch}
                    $\rightarrow$ \eng{decode}
              \item \eng{fetch} $\rightarrow$ \eng{decode}
                    $\rightarrow$ \eng{execute}
              \item \eng{fetch} $\rightarrow$ \eng{execute}
                    $\rightarrow$ \eng{decode}
          \end{enumerate}

    \item Nilai \dec{-10} dalam 8-bit \eng{two's complement} adalah
          \dots
          \begin{enumerate}
              \item \hex{0A} \item \hex{0F}
              \item \hex{F5} \item \hex{F6}
          \end{enumerate}

    \item Hasil \asm{ADD AL, 01h} dengan \reg{AL} = \hex{7F} adalah
          \dots
          \begin{enumerate}
              \item \reg{AL} = \hex{80}, \flag{OF} = 0,
                    \flag{CF} = 0
              \item \reg{AL} = \hex{80}, \flag{OF} = 1,
                    \flag{CF} = 0
              \item \reg{AL} = \hex{80}, \flag{OF} = 1,
                    \flag{CF} = 1
              \item \reg{AL} = \hex{00}, \flag{OF} = 0,
                    \flag{CF} = 1
          \end{enumerate}

    \item Sandian ASCII untuk karakter \code{'A'} adalah \dots
          \begin{enumerate}
              \item \hex{0A} \item \hex{41}
              \item \hex{61} \item \hex{45}
          \end{enumerate}
\end{enumerate}

\subsection{Soal Uraian}

\begin{enumerate}
    \item Jelaskan model mesin penyimpan program beserta keempat
          komponennya. Kemudian jelaskan satu konsekuensi penting
          dari prinsip \eng{stored-program} bagi cara program dimuat
          dan dijalankan.

    \item Jelaskan perbedaan antara \eng{carry} dan \eng{overflow}.
          Sertakan satu contoh operasi yang menghasilkan \flag{CF} = 1
          tetapi \flag{OF} = 0, dan satu contoh yang menghasilkan
          \flag{CF} = 0 tetapi \flag{OF} = 1.

    \item Sebuah variabel berisi \hex{32}. Jelaskan dua tafsiran yang
          mungkin atas isi variabel tersebut, dan sebutkan konteks
          yang menentukan tafsiran mana yang benar.

    \item Berdasarkan Gambar~\ref{fig:peta-konsep}, susun rencana
          belajar Anda untuk enam minggu pertama. Jelaskan mengapa
          urutan bab tidak dapat ditukar sembarangan.
\end{enumerate}

\subsection{Kisi-kisi Penilaian}

\begin{table}[H]
    \centering
    \caption{Kisi-kisi evaluasi Bab~\ref{chap:landasan-teori}.}
    \label{tab:kisi-bab2}
    \small
    \begin{tabular}{@{}p{5.4cm}p{5.2cm}c@{}}
        \toprule
        \textbf{Aspek yang diukur} & \textbf{Bentuk soal} &
        \textbf{Bobot} \\
        \midrule
        Konsep model mesin dan siklus instruksi & Pilihan ganda 1, 2;
        Uraian 1 & 25\% \\
        \addlinespace
        Representasi bilangan bertanda & Pilihan ganda 3; Uraian 2 &
        25\% \\
        \addlinespace
        Pembedaan \eng{carry} dan \eng{overflow} & Pilihan ganda 4;
        Uraian 2 & 20\% \\
        \addlinespace
        Penyandian data dan tafsirannya & Pilihan ganda 5; Uraian 3 &
        15\% \\
        \addlinespace
        Kemampuan membaca peta konsep buku & Uraian 4 & 15\% \\
        \bottomrule
    \end{tabular}
\end{table}

\begin{checklistbox}
\textbf{Checklist Pencapaian Fondasi Bab~\ref{chap:landasan-teori}}

Mahasiswa dinyatakan siap memasuki bab materi apabila:
\begin{itemize}
    \item[\checked] dapat menjelaskan model von Neumann dan siklus
          instruksi;
    \item[\checked] dapat mengonversi bilangan biner, heksadesimal,
          dan desimal dengan benar tanpa alat bantu;
    \item[\checked] dapat menyatakan bilangan negatif dalam
          \eng{two's complement} dan memverifikasi hasilnya;
    \item[\checked] dapat membedakan \eng{carry} dan \eng{overflow}
          beserta flag yang menandainya;
    \item[\checked] dapat menjelaskan perbedaan nilai biner dan
          sandian karakter;
    \item[\checked] dapat menjelaskan urutan ketergantungan antarbab
          melalui peta konsep.
\end{itemize}

\vspace{0.4em}
Seluruh butir di atas bersifat prasyarat untuk
Bab~\ref{chap:pengantar-assembly}.
\end{checklistbox}


\end{document}
